% !PRINTED 455
\weberpage{455}

%<<CONT
\gesperrt{bare, die der ersten Art als unzerlegbare Functionen} bezeichnet werden. Ist z.\,B. $\Omega$ der Körper der rationalen Zahlen, so ist $x^2 - y^2 = (x - y)(x + y)$ reducibel in $\Omega$. Die Function $x^2 - 2y^2$ ist zwar in die Factoren $x + y\sqrt{2}, x - y\sqrt{2}$ zerlegbar, aber diese Factoren sind nicht Functionen in $\Omega$. Die Function $x^2 + y^2 + 1$ endlich ist überhaupt nicht zerlegbar. Eine in einem Körper $\Omega$ irreducible, aber zerlegbare Function wird in einem erweiterten Körper reducibel.

Die fundamentalen Sätze über irreducible Functionen haben wir in \S.\,51 kennen gelernt. Bei jenen Ausführungen war zwar zunächst, dem dortigen Standpunkte entsprechend, der Körper aller Zahlen als Rationalitätsbereich vorausgesetzt, in dem der Unterschied zwischen unzerlegbaren und irreducibeln Functionen wegfällt. Benutzt ist aber immer nur die Reducibilität oder Irreducibilität einer Function, und die besondere Natur des Körpers $\Omega$ kommt nicht in Betracht. So ist also darin der Beweis des Satzes enthalten:

\gesperrt{IV. Eine reducible Function in $\Omega$ kann nur auf eine Art in irreducible Factoren zerlegt werden, wenn man rationale Functionen, die sich nur durch einen constanten Factor unterscheiden, als nicht verschieden betrachtet.}

\begin{center}
§. 142.\\
\gesperrt{Algebraische Körper.}
\end{center}

Ist $\Omega$ ein beliebiger Körper, und $F(x)$ eine Function in $\Omega$, so heisst die Gleichung
\weq{1}{F(x) = 0}
eine \gesperrt{Gleichung in} $\Omega$. Diese Gleichung heisst reducibel oder irreducibel, je nachdem die Function $F(x)$ reducibel oder irreducibel ist. Durch Adjunction einer Wurzel $\alpha$ einer solchen Gleichung zu $\Omega$ entsteht (wenn $\alpha$ nicht selbst schon zu $\Omega$ gehört) ein neuer Körper $\Omega'$, den wir einen \gesperrt{algebraischen Körper} \glqq\gesperrt{über}\grqq{} $\Omega$ oder auch, wenn Zweifel über die Bedeutung ausgeschlossen sind, kurz, einen \gesperrt{algebraischen Körper} nennen wollen. Wir brauchen für einen solchen Körper das Zeichen
\[ \Omega(\alpha). \]
